%% plantilla Beamer de  https://code.google.com/p/pyprop/ Tore Birkeland


\newcommand{\seccion}[1] {
\section{#1}

\begin{frame} % frame
        \begin{center}
                \Huge
                \textcolor{lbluegreen}{#1}
        \end{center}
\end{frame}
}

\documentclass{beamer}
\usepackage{pgf,pgfpages}
\usepackage{graphicx}
\usepackage{units}
\usepackage[utf8]{inputenc}
\usepackage{mathtools}
\usepackage{amsmath}
\setbeamertemplate{title page}
{
  \vbox{}
  \vfill
 \begin{raggedright}
      \vskip0.25em%
      {\usebeamercolor[fg]{titlegraphic}\inserttitlegraphic\par}        
      \vskip13em%
    \begin{beamercolorbox}[sep=4pt,center]{title}
      \usebeamerfont{title}\inserttitle\par%
      \ifx\insertsubtitle\@empty%
      \else%
        \vskip0.25em%
        {\usebeamerfont{subtitle}\usebeamercolor[fg]{subtitle}\insertsubtitle\par}%
      \fi%    
    \end{beamercolorbox}%
    \vskip1em\par
    \begin{beamercolorbox}[sep=4pt,center]{author}
      \usebeamerfont{author}\insertauthor
    \end{beamercolorbox}
    \begin{beamercolorbox}[sep=4pt,center]{institute}
      \usebeamerfont{institute}\insertinstitute
    \end{beamercolorbox}
    
%\vspace{-20pt} %<-- right here
    \begin{beamercolorbox}[sep=8pt,center]{date}
      \usebeamerfont{date}\insertdate
    \end{beamercolorbox}\vskip0.5em
   \end{raggedright}
  \vfill
}


\mode<presentation>
{
  \usetheme{ift}
  \setbeamercovered{transparent}
  \setbeamertemplate{items}[square]
}

\usefonttheme[onlymath]{serif}
\setbeamerfont{frametitle}{size=\LARGE,series=\bfseries}

\definecolor{uibred}{RGB}{170, 0, 0}
\definecolor{uibblue}{RGB}{0, 84, 115}
\definecolor{uibgreen}{RGB}{119, 175, 0}
%\definecolor{uibgreen}{RGB}{50, 105, 0}
\definecolor{uiborange}{RGB}{217, 89, 0}


\beamertemplatenavigationsymbolsempty


%====================================================
%-------------Macro definitions go here--------------
%====================================================

%
% Differentials
%
\newcommand{\tdiff}[2]{\ensuremath{\frac{\mathrm{d}#2}}{\mathrm{d}{#1}}}
\newcommand{\tdifforder}[3]{\ensuremath{\frac{\mathrm{d}^{#2}#3}{\mathrm{d}{#1}^{#2}}}}
\newcommand{\pdiff}[2]{\ensuremath{\frac{\partial#2 }{\partial#1}}}
\newcommand{\pdifforder}[3]{\ensuremath{\frac{\partial^{#2}#3}{\partial{#1}^{#2}}}}

% bracket
\newcommand{\bra}[1]{\ensuremath{\left<#1\right|}}
\newcommand{\ket}[1]{\ensuremath{\left|#1\right>}}
\newcommand{\bracket}[2]{\ensuremath{ \left\langle #1 | #2 \right\rangle}}
\newcommand{\matelem}[3]{\ensuremath{ \left\langle #1 | #2 | #3 \right\rangle}}
\newcommand{\matr}[1]{\ensuremath{\mathbf{#1}}}
\newcommand{\vect}[1]{\ensuremath{\mathbf{#1}}}
\newcommand{\expectationvalue}[1]{\ensuremath{\left\langle #1 \right\rangle}}

\renewcommand{\imath}{\ensuremath{\mathrm{i}}}

%
% Linear algebra
%
\renewcommand{\vec}[1]{\ensuremath{\mathbf{#1}}}
\newcommand{\mat}[1]{\ensuremath{\mathbf{#1}}}
\newcommand{\tildemat}[1]{\ensuremath{\widetilde{\mat{#1}}}}

% Numerical analysis
\newcommand{\bigo}{\ensuremath{\mathcal{O}}}
\renewcommand{\Re}{\ensuremath{\mathrm{Re}}}
\renewcommand{\Im}{\ensuremath{\mathrm{Im}}}
\newcommand{\mathcol}[2]{{\color{#1}#2}}
\newcommand{\red}[1]{\mathcol{uibred}{#1}}
\newcommand{\blue}[1]{\mathcol{uibblue}{#1}}
\newcommand{\green}[1]{\mathcol{uibgreen}{#1}}
\newcommand{\orange}[1]{\mathcol{uiborange}{#1}}
%
% Misc macros
%
\newcommand{\eref}[1]{~(\ref{#1})}
\renewcommand{\equiv}[0]{\ensuremath{:=}}
\newcommand{\etal}{\textit{et al. }}
\newcommand{\paperheader}[2]{\noindent\textbf{Paper #1}: \textit{#2}\\}
\newcommand{\paperitem}[3]{\noindent\textbf{Paper #1}: \textit{#2}\vspace{1em}\\\noindent #3\vspace{2em}}
\newcommand{\tfinal}{\ensuremath{T_{\text{f}}}}
\newcommand{\papernum}[1]{\textbf{#1}}
%\newcommand{\note}[1]{\colorbox{yellow}{#1}}
\newcommand{\paperref}[1]{Paper~\textbf{#1}}

%
% Code
%
\newcommand{\inlinename}[1]{\lstinline[basicstyle=\ttfamily,language=bash]{#1}}




%\includeonlyframes{current}


\defbeamertemplate{enumerate item}{mycircle}
{
  %\usebeamerfont*{item projected}%
  %\usebeamercolor[bg]{item projected}%
  \begin{pgfpicture}{0ex}{0ex}{1.5ex}{0ex}
	%\pgfcircle[fill]{\pgfpoint{0pt}{.75ex}}{1.25ex}
    \pgfbox[center,base]{\color{uibblue}\insertenumlabel.}
  \end{pgfpicture}%
}
[action]
{\setbeamerfont{item projected}{size=\scriptsize}}
\setbeamertemplate{enumerate item}[mycircle]




\title{The solution for the construction topology problem for rural wireless networks}
\author{Milton J. Ríos Rivera \and Leonardo Rodríguez Mújica }
\institute{Universidad de Cundinamarca}
\date{Abril 2015}

\begin{document}


\setbeamertemplate{background}
 {\includegraphics[width=\paperwidth,height=\paperheight]{portada.pdf}}
\setbeamertemplate{footline}[default]

\begin{frame}
  \titlepage
  \vspace{5cm}
\end{frame}

%
% Set the background for the rest of the slides.
% Insert infoline at the end
%
\setbeamertemplate{background}
 {\includegraphics[width=\paperwidth,height=\paperheight]{fondo}}
\setbeamertemplate{footline}[ifttheme]

%--------------------------------------------------------------------
%                          Introduction
%--------------------------------------------------------------------

\seccion{Introducción}


%--------------------------------------------------------------------
%                          Outline
%--------------------------------------------------------------------

\subsection{Motivaciones}

\begin{frame}
	\frametitle{¿Por qué?}
	
	\begin{columns}
		\column{0.6\textwidth}
			\parbox[c][0.8\textheight]{0.8\textwidth}
			{
				\includegraphics<1>[width=0.5\textwidth]{figurer/udec-logo.jpg}
%% http://upload.wikimedia.org/wikipedia/commons/7/7b/Mapa_Provincia_del_Sumapaz.JPG
				\includegraphics<2>[width=0.7\textwidth]{figurer/mapa-sumapaz.jpg}
%% http://www.mintic.gov.co/portal/vivedigital/612/w3-propertyvalue-642.html
				\includegraphics<3>[width=\textwidth]{figurer/tics.png}
%% http://www.mintic.gov.co/portal/vivedigital/612/articles-1510_recurso_1.pdf pág 31
				\includegraphics<4>[width=\textwidth]{figurer/estudiante.jpg}
				\includegraphics<5>[width=0.5\textwidth]{figurer/torres.jpg}
			}
		
		
		\column{0.4\textwidth}
			\begin{tabular}{l}
				 \onslide<1->{La institución} \\
				 \onslide<2->{La región del Sumapaz} \\
				 \onslide<3->{Brecha digital} \\
				 \onslide<4->{I. E. Rurales} \\
				 \onslide<5->{Infraestructura} \\
			\end{tabular}
	\end{columns}

\end{frame}

\seccion{ Problema de construcción  de la topología}


\subsection{Requerimientos}
\begin{frame}
	\frametitle{Requerimientos del problema}

	\begin{columns}
		\column{0.5\textwidth}
			\parbox[c][0.8\textheight]{0.9\textwidth}
			{
				\includegraphics<1>[width=\textwidth]{figurer/nodos.pdf}
				\includegraphics<2>[width=\textwidth]{figurer/tipos-de-torres.jpg}
				\includegraphics<3>[width=\textwidth]{figurer/torres.pdf}
				\includegraphics<4>[width=\textwidth]{figurer/fresnel.pdf}
			}
		
		
		\column{0.5\textwidth}
			\begin{tabular}{l}
				 \onslide<1->{Conectividad} \\
				 \onslide<2->{Altura máxima de las torres} \\
				 \onslide<3->{Naturaleza de la función de costos} \\
				 \onslide<4->{Pérdidas por espacio libre y} \\
				 \onslide<4->{línea de vista} \\
			\end{tabular}
	\end{columns}
\end{frame}



\subsection{Complejidad de Calculo}

\begin{frame} % frame
	\frametitle{Complejidad de Calculo}


    Los autores de del algorimo: Panigrahi-Duttat-Naiswal-Naidu-Rastogi
    determinan que el algoritmo óptimo
    es  {\em NP-hard}, con una reducción del
    conjunto de vértices conectados no
    mejor que un factor logarítmico $O(\log n)$, sin embargo, demuestran
    que por medio de
    una aproximación se consigue un comportamiento de reducción constante a
    través de un algoritmo voraz llamado TC-ALGO 

\end{frame}

%\begin{itemize}
%		\item 
%		\item 
%		\item 
%		\item 
%		\item 
%	\end{itemize}


\seccion{Ecuación de Costos}

\subsection{Ecuación de costos: aplicada a la región}
       
\begin{frame}
	\frametitle{Costos en Colombia}

        \parbox[c][2.5cm]{\textwidth}
        {
             $$
                        c(h)=\left\{
                        \begin{array}{ll}
                            250000 & si \,\,0 < \leq h \leq 20m\\
                            550000h+B & si \,\, 20 < h \leq 60 m
                        \end{array}
                        \right .
                    $$
        }

        \parbox[c][4cm]{\textwidth}
        {
                \begin{center}
                        \includegraphics<1>[width=0.7\textwidth]{figurer/torres.pdf}
                \end{center}
        }
\end{frame}

       
\seccion{Descripción del algoritmo} 



\begin{frame}
        \frametitle{Algoritmo TC-ALGO}

        \begin{columns}
                \column{0.5\textwidth}
                        \parbox[c][0.8\textheight]{0.9\textwidth}
                        {
                                \includegraphics<1->[width=\textwidth]{figurer/star_tc_algo.pdf}
                        }
                
                 
                \column{0.5\textwidth}
                                \onslide<1->{$G(V,E)$ determina la estructura de la entrada de la red, }   
                                \onslide<2->{donde $V$ es el conjunto de nodos, }  
                                \onslide<3->{$E$ es el conjunto de aristas} 
                                \onslide<4->{y $|V|=n$ número de nodos.} \\
                                \onslide<5->{Un subgrafo denominado
                                $COVER(h)$ es el conjunto de aristas que
                            son establecidas por la función de altura $h$. }
                                \onslide<6->{Al inicio $h=0$, $COVER(h)
                                =\emptyset$ y $COMP(h)=n$.}
                                \onslide<7->{El algoritmo finaliza cuando $COMP(h)=1$}
        \end{columns}
\end{frame}

\subsection{El problema}

\begin{frame}
    \frametitle{El problema a resolver}
    {\Large
El problema es encontrar una función de altura que represente el menor
    costo total posible, para un subgrafo conectado y expandido.
$$min \left( \sum_{u \in V} c[h(u)] \right) $$}
\end{frame}    

\seccion{Desarrollo de la aplicación}

\subsection{Topología inicial}

\begin{frame} % frame
	\frametitle{Topología de la red}
	

   	\begin{columns}
		\column{0.5\textwidth}
			\parbox[c][0.8\textheight]{0.9\textwidth}
			{
				\includegraphics<1>[width=0.9\textwidth]{figurer/topologia.pdf}
				\includegraphics<2>[width=\textwidth]{figurer/red-veredas-simple.pdf}
				\includegraphics<3>[width=0.9\textwidth]{figurer/veredas.pdf}
				\includegraphics<4>[width=\textwidth]{figurer/topo_inicial.pdf}
			}
		
		
		\column{0.5\textwidth}
			\begin{tabular}{l}
				 \onslide<1->{Conectividad VPN} \\
				 \onslide<2->{Red en veredas} \\
				 \onslide<3->{Ubicación Geográfica} \\
				 \onslide<4->{Mapas SRTM} \\
			\end{tabular}
	\end{columns}
\end{frame}

\subsection{Características del diseño de la topología}

\begin{frame} % frame
	\frametitle{Grafo inicial}
	\begin{itemize}
		\item <1-> El grafo de entrada
 $G_i (V_i , E_i)$ se establece completamente conectado 
		\item<2->  El costo de este enlace provienen de su altura
		\item <3->Decartar los enlaces que no
cumplen el valor de altura máxima $h_{max}$ 
		\item <4->$h_o (u, v)$ es la función de obstrucciones en el trayecto del enlace
la cual describe el perfil del terreno entre los nodos
$u$ y $v$
	\end{itemize}
\end{frame}

\subsection{Algoritmo de optimización TC-ALGO}

\begin{frame} % frame
	\parbox[c][5cm]{\textwidth}
	{
        \frametitle{Resultados}
		\begin{center}
			\includegraphics<1>[width=8cm]{figurer/grafo-inicial.pdf}
			\includegraphics<2>[width=8cm]{figurer/resul-tcalgo2.pdf}
			\includegraphics<3>[width=7cm]{figurer/cabrera.png}
			\includegraphics<4>[width=7cm]{figurer/ingcabreranl.png}
			\includegraphics<5>[width=7cm]{figurer/covercabrera.png}
		\end{center}
	}

    \begin{center}   
     \only<1> {Enlaces que se pierden debido a la topografía del
    terreno}
    \only <2> {Resultado de costos y de densidad de los grafos
        resultantes}
    \only<3> {Grafo de entrada $G$ de Cabrera totalmente
        conectado}
    \only<4> {Grafo inicial $G_i$ de Cabrera luego de
            seleccionar los nodos con restricción de altura} 
    \only<5> {Grafo de salida de Cabrera con los enlaces
            cubiertos}
    \end{center}        
\end{frame}

\seccion{Trabajo futuro} 

\begin{frame} % frame
	\frametitle{Trabajo futuro}
     {\bf Problema de conectividad}\\
      Debido a que en el algoritmo se descartan nodos que se esperaban
       inicialmente estuvieran conectados a la topologia de red final, esto
       se 
        soluciona agragando nodos a la red en puntos geograficos
        especificos que
         agreguen dicho beneficio de conectividad. Dado esto se pretende:

	\begin{itemize}
		\item <1-> Encontrar puntos en el terreno que
permitan conectar los nodos que en un principio resultaron inalcanzables
bajo las consideraciones de este trabajo.
		\item <2->Determinar la eficiencia de dicho algoritmo 
	\end{itemize}
\end{frame}

\subsection{Créditos} 

\begin{frame} % frame
	\frametitle{Origen de las imágenes}
    {\bf Wikimedia Commons}

\begin{itemize}
        \item Mapa del Sumapaz - File:Mapa\_Provincia\_del\_Sumapaz.JPG
    \end{itemize}

{\bf MinTic }

\begin{itemize}
        \item Joven con computador - articles-1510\_recursoi\_1.pdf
        \item Tabla hitos y metas - http://www.mintic.gov.co/portal/vivedigital/612/w3-propertyvalue-642.html
\end{itemize}

{\bf Plantilla Beamer}

\begin{itemize}
    \item  Tore Birkeland. https://code.google.com/p/pyprop/
\end{itemize}

\end{frame}    

\section{ Preguntas}
\subsection{Preguntas}
\begin{frame}
				\includegraphics[width=0.8\textwidth]{preguntas.pdf}

\end{frame}

\end{document}
